\chapter{Peubah Acak Diskret}\label{ch:b2:randomvar}

\omterm{def:b2:randomvar:law}{Peubah acak} menata perhitungan peluang di sekitar
fungsi alih-alih \omterm{def:b2:proba:space}{kejadian}. Pada ruang \omterm{def:b2:structures:countable}{terbilang} teorinya
digerakkan \omterm{def:b2:series:summable}{keluarga terjumlahkan} pada \cref{ch:b2:series}:
sebab \omterm{def:b2:randomvar:expectation}{nilai harapan} adalah jumlah sebuah keluarga yang diindeks \omterm{def:b2:proba:space}{ruang sampelnya},
dan semua sifatnya --- kelinearan, pemindahan, rumus hasil kali
bagi peubah yang \omterm{def:b2:proba:independence}{saling bebas} --- merupakan teorema tentang keluarga
yang \omterm{def:b2:series:summable}{terjumlahkan}. Bab ini membuktikan ketaksamaan kunci Markov,
Chebyshev, Cauchy--Schwarz dan Jensen, lalu berakhir dengan \omterm{def:b2:randomvar:law}{hukum}
klasiknya beserta \omterm{def:b2:randomvar:law}{hukum} bilangan besar yang lemah, yang buktinya dua baris
begitu Chebyshev tersedia.

\section{Peubah acak dan hukumnya}

\begin{definition}[Peubah acak diskret; hukumnya]\label{def:b2:randomvar:law}
Misalkan $(\Omega, \P)$ \omterm{def:b2:proba:space}{ruang peluang} yang \omterm{def:b2:structures:countable}{terbilang}. Sebuah
\emph{peubah acak}\index{peubah acak} adalah pemetaan $X \colon \Omega \to E$ (dengan $E$ sembarang
himpunan; dan disebut \emph{real} bila $E = \R$). Adapun \emph{hukum}\index{hukum!peubah acak}
(atau \emph{distribusi}\index{distribusi}) peubahnya adalah \omterm{def:b2:proba:space}{ukuran peluang} $\P_X$ pada
\omterm{def:b2:structures:countable}{himpunan terbilang} $X(\Omega)$ yang didefinisikan lewat
\[
\P_X(\{x\}) = \P(X = x)
= \P\bigl(\{\omega : X(\omega) = x\}\bigr) .
\]
\end{definition}

\begin{example}[Hukum klasiknya]\label{ex:b2:randomvar:laws}
\leavevmode
\begin{itemize}
\item \emph{Bernoulli} $\mathcal{B}(p)$: dengan $X \in \{0, 1\}$ dan $\P(X =
1) = p$. Yakni indikator sebuah \omterm{def:b2:proba:space}{kejadian}.
\item \emph{Binomial} $\mathcal{B}(n, p)$: dengan $\P(X = k) =
\binom nk p^k(1-p)^{n-k}$ untuk $0 \leq k \leq n$: yakni banyaknya keberhasilan
pada $n$ percobaan Bernoulli yang \omterm{def:b2:proba:independence}{saling bebas} (jilid Kelas 10--12; yang dibuktikan ulang
di bawah lewat jumlah peubah yang \omterm{def:b2:proba:independence}{saling bebas}).
\item \emph{Geometri} $\mathcal{G}(p)$: dengan $\P(X = k) =
(1-p)^{k-1}p$ untuk $k \in \N^*$: yakni pangkat keberhasilan pertamanya
(\cref{ex:b2:proba:geometric}).
\item \emph{Poisson} $\mathcal{P}(\lambda)$: dengan $\P(X = k) =
e^{-\lambda}\frac{\lambda^k}{k!}$ untuk $k \in \N$ --- yakni sebuah \omterm{def:b2:proba:space}{ukuran
peluang} berkat deret eksponensialnya. Inilah \omterm{def:b2:randomvar:law}{hukum} \omterm{def:b2:proba:space}{kejadian} langka
(\cref{ch:b2:genfun}).
\end{itemize}
\end{example}

\begin{remark}[Hukum mana memodelkan apa]
Keempat \omterm{def:b2:randomvar:law}{hukumnya} menjawab empat pertanyaan purba: Bernoulli,
``apakah ia terjadi?''; binomial, ``berapa kali dalam $n$
percobaan?''; geometri, ``berapa lama sampai kali pertamanya?'';
dan Poisson, ``berapa banyak \omterm{def:b2:proba:space}{kejadian} pada laju tertentu, ketika percobaannya
banyak dan masing-masing tak mungkin?''. Jadi mengenali pertanyaannya
adalah sembilan persepuluh pemodelannya: sebab jumlah indikator menunjuk ke
binomialnya, waktu tunggu ke geometrinya, dan cacah \omterm{def:b2:proba:space}{kejadian}
langka ke Poissonnya --- dengan peralihan dari binomial ke
Poisson dipersis oleh \omterm{def:b2:randomvar:law}{hukum} \omterm{def:b2:proba:space}{kejadian} langka pada
\cref{ch:b2:genfun}.
\end{remark}

\begin{proposition}[Sifat tanpa ingatan pada hukum
geometri]\label{prop:b2:randomvar:memoryless}
Bila $X \sim \mathcal{G}(p)$, maka untuk setiap $m, n \in \N$:
\[
\P(X > m + n \mid X > m) = \P(X > n) ,
\]
dan \omterm{def:b2:randomvar:law}{hukum} geometri adalah satu-satunya \omterm{def:b2:randomvar:law}{hukum} pada $\N^*$ yang bersifat
demikian.
\end{proposition}

\begin{proof}
Menjumlahkan bobot geometrinya memberi $\P(X > n) = (1-p)^n$. Karena itu
\[
\P(X > m + n \mid X > m)
= \frac{\P(X > m + n)}{\P(X > m)}
= \frac{(1-p)^{m+n}}{(1-p)^m} = (1-p)^n = \P(X > n).
\]
Sebaliknya, bila $G(n) = \P(X > n)$ memenuhi $G(m + n) =
G(m)G(n)$ dengan $G(0) = 1$, maka $G(n) = G(1)^n$ secara induktif; lalu $q =
G(1) \in \intco{0}{1}$, dan $q = 0$ atau \omterm{def:b2:randomvar:law}{hukumnya} adalah
$\mathcal{G}(1 - q)$: sebab $\P(X = k) = G(k-1) - G(k) =
q^{k-1}(1 - q)$.
\end{proof}

\begin{example}[Tak ada angka yang pernah ``jatuh tempo'']
Lemparkan sebuah dadu sambil menunggu enam: maka waktu tunggunya $X \sim
\mathcal G(1/6)$. Adapun sifat tanpa ingatannya mengatakan bahwa sesudah $10$
lemparan yang sia-sia, tunggu \emph{sisanya} $X - 10$, jika diketahui
$X > 10$, kembali berhukum $\mathcal G(1/6)$: jadi \omterm{def:b2:randomvar:expectation}{nilai harapan}
tunggu bersyaratnya masih $6$ lemparan, persis seperti pada awalnya.
Dadunya tak mengingat, dan tak ada enam yang pernah ``jatuh tempo'' --- jadi
kekeliruan penjudi adalah kepercayaan bahwa \omterm{def:b2:randomvar:law}{hukum} bersyaratnya
sepatutnya bergeser. Sebaliknya, paruh ketunggalan proposisinya
mengatakan bahwa ketidakpedulian ini \emph{mencirikan} waktu tunggu
geometri: sebab sembarang waktu tunggu yang ramalannya tak pernah diperbarui
bersifat geometri. Adapun antrean dan umur pakai yang nyata biasanya memperbarui,
dan persis begitulah kita mendeteksi bahwa keduanya tak geometri.
\end{example}

\section{Nilai harapan}

\begin{definition}[Nilai harapan]\label{def:b2:randomvar:expectation}
Sebuah \omterm{def:b2:randomvar:law}{peubah acak} real $X$ pada $(\Omega, \P)$ disebut \emph{punya
nilai harapan} bila keluarga
$\bigl(X(\omega)\,\P(\{\omega\})\bigr)_{\omega \in \Omega}$ bersifat
\omterm{def:b2:series:summable}{terjumlahkan} (\cref{ch:b2:series}); dan \emph{nilai harapannya}\index{nilai harapan} lalu
\[
\E(X) = \sum_{\omega \in \Omega} X(\omega)\,\P(\{\omega\}) .
\]
\end{definition}

\begin{theorem}[Teorema pemindahan]\label{thm:b2:randomvar:transfer}
Peubah $X$ punya \omterm{def:b2:randomvar:expectation}{nilai harapan} bila dan hanya bila keluarga $\bigl(x\,\P(X =
x)\bigr)_{x \in X(\Omega)}$ \omterm{def:b2:series:summable}{terjumlahkan}, dan lalu
\[
\E(X) = \sum_{x \in X(\Omega)} x\,\P(X = x) .
\]
Lebih umum, untuk $f \colon X(\Omega) \to \R$, peubah
$f(X)$ punya \omterm{def:b2:randomvar:expectation}{nilai harapan} bila dan hanya bila $\sum_x \abs{f(x)}\,\P(X = x) <
\infty$, dan lalu $\E(f(X)) = \sum_x f(x)\,\P(X = x)$.
\end{theorem}

\begin{proof}
Partisikanlah $\Omega$ menjadi himpunan arasnya $\Omega_x = \{X = x\}$ dengan $x
\in X(\Omega)$. Menurut teorema penjumlahan lewat paket bagi keluarga
yang \omterm{def:b2:series:summable}{terjumlahkan} (\cref{ch:b2:series}), keluarga
$(X(\omega)\P(\{\omega\}))_\omega$ \omterm{def:b2:series:summable}{terjumlahkan} bila dan hanya bila tiap paketnya
\omterm{def:b2:series:summable}{terjumlahkan} (yang otomatis: sebab $\sum_{\omega \in \Omega_x}\abs{x}\P(\{\omega\}) =
\abs x\,\P(X = x)$) \emph{dan} keluarga jumlah paketnya
$\bigl(x\,\P(X = x)\bigr)_x$ \omterm{def:b2:series:summable}{terjumlahkan} --- dan lalu jumlah
totalnya sepakat. Adapun untuk $f(X)$: terapkanlah pernyataan yang terbukti itu pada peubah $Y =
f \circ X$, yang himpunan arasnya $\{Y = y\} =
\bigsqcup_{x : f(x) = y}\{X = x\}$; lalu penjumlahan lewat paket yang
kedua mengubah $\sum_y y\,\P(Y = y)$ menjadi $\sum_x
f(x)\,\P(X = x)$, dengan paketnya kini mengelompokkan nilai $x$
menurut petanya $f(x)$, dan keterjumlahan \omterm{def:b2:series:def}{mutlak} satu
keluarganya setara dengan keterjumlahan yang lain.
\end{proof}

\begin{theorem}[Sifat nilai harapan]\label{thm:b2:randomvar:linearity}
Pada himpunan \omterm{def:b2:randomvar:law}{peubah acak} yang punya \omterm{def:b2:randomvar:expectation}{nilai harapan}:
\begin{enumerate}
\item (Kelinearan) $\E(aX + bY) = a\,\E(X) + b\,\E(Y)$.
\item (Kepositifan dan kemonotonan) $X \geq 0 \Rightarrow \E(X)
\geq 0$; $X \leq Y \Rightarrow \E(X) \leq \E(Y)$; dan
$\abs{\E(X)} \leq \E(\abs X)$.
\item (Dominasi) Bila $\abs X \leq Z$ dan $Z$ punya \omterm{def:b2:randomvar:expectation}{nilai harapan},
maka $X$ juga.
\end{enumerate}
\end{theorem}

\begin{proof}
Semuanya merupakan sifat jumlah keluarga yang \omterm{def:b2:series:summable}{terjumlahkan}
(\cref{ch:b2:series}): yakni kelinearan jumlahnya, kepositifan suku demi
suku, dan kriteria dominasi bagi keterjumlahannya. (Perhatikan bahwa
kelinearannya seketika pada \emph{definisi}nya atas $\Omega$,
sedangkan ia akan janggal pada rumus pemindahannya --- yakni satu keuntungan
mendefinisikan $\E$ di hulu.)
\end{proof}

\begin{example}\label{ex:b2:randomvar:expectations}
Untuk $X \sim \mathcal{B}(n, p)$: dengan menulis $X = X_1 + \dots + X_n$ sebagai
jumlah indikator Bernoulli lalu memakai kelinearannya, $\E(X) = np$ ---
tanpa perlu koefisien binomial. Untuk $X \sim \mathcal{G}(p)$: $\E(X) =
\sum_{k\geq1}k(1-p)^{k-1}p = p\cdot\frac{1}{(1 - (1-p))^2} =
\frac1p$, lewat menurunkan deret geometrinya di dalam cakramnya
(\cref{ch:b2:powerseries}). Untuk $X \sim \mathcal{P}(\lambda)$: $\E(X)
= \sum_{k\geq1}k e^{-\lambda}\frac{\lambda^k}{k!} = \lambda
e^{-\lambda}\sum_{j\geq0}\frac{\lambda^j}{j!} = \lambda$.
\end{example}

\begin{example}[Pemindahan beraksi]\label{ex:b2:randomvar:transferex}
Untuk $X \sim \mathcal P(\lambda)$, hitunglah
$\E\bigl(\frac1{1+X}\bigr)$ --- sebab \omterm{def:b2:randomvar:law}{hukum} $\frac1{1+X}$
sendiri janggal, tetapi pemindahannya tak pernah memintanya:
\[
\E\Bigl(\frac1{1+X}\Bigr)
= \sum_{k\geq0}\frac{1}{k+1}\,\eu^{-\lambda}
\frac{\lambda^k}{k!}
= \frac{\eu^{-\lambda}}{\lambda}\sum_{k\geq0}
\frac{\lambda^{k+1}}{(k+1)!}
= \frac{\eu^{-\lambda}}{\lambda}\bigl(\eu^\lambda - 1\bigr)
= \frac{1 - \eu^{-\lambda}}{\lambda} .
\]
Ada dua pelajaran. Secara hitungan: mengenali deret eksponensial
yang tergeser adalah seluruh kerjanya --- sebab pemindahan menyusutkan
\omterm{def:b2:randomvar:expectation}{nilai harapan} $f(X)$ menjadi manipulasi deret. Secara struktural:
nilai colok naifnya adalah $\frac1{1 + \E X} =
\frac1{1 + \lambda}$, sedangkan jawaban benarnya lebih besar,
\[
\frac{1 - \eu^{-\lambda}}{\lambda} \geq
\frac{1}{1 + \lambda},
\]
persis seperti dituntut ketaksamaan Jensen bagi fungsi \omterm{def:b2:affine:convex}{cembung} $t
\mapsto \frac1{1+t}$. Jadi \omterm{def:b2:randomvar:expectation}{nilai harapan} peta
\omterm{def:b2:affine:convex}{cembung} duduk di atas nilai colok naifnya, dan pemindahan
ditambah pemeriksaan deret membuat ketaksamaan abstraknya menjadi konkret.
\end{example}

\begin{theorem}[Kebebasan dan hasil kali]\label{thm:b2:randomvar:product}
\omterm{def:b2:randomvar:law}{Peubah acak} $X, Y$ disebut \emph{\omterm{def:b2:proba:independence}{saling bebas}} bila $\P(X = x, Y =
y) = \P(X = x)\P(Y = y)$ untuk setiap $x, y$ --- setara dengan itu, bila
\omterm{def:b2:proba:space}{kejadian} $\{X \in A\}$ dan $\{Y \in B\}$ \omterm{def:b2:proba:independence}{saling bebas} untuk setiap $A,
B$. Bila $X$ dan $Y$ merupakan peubah real yang \omterm{def:b2:proba:independence}{saling bebas} dan punya
\omterm{def:b2:randomvar:expectation}{nilai harapan}, maka $XY$ punya \omterm{def:b2:randomvar:expectation}{nilai harapan} dan
\[
\E(XY) = \E(X)\,\E(Y) .
\]
\end{theorem}

\begin{proof}
Kesetaraan kedua perumusannya menyusul dengan menjumlahkan
kesamaan \omterm{def:b2:funcseq:def}{titik demi titiknya} atas $(x, y) \in A \times B$
(lewat keaditifan-$\sigma$ dua kali). Adapun untuk hasil kalinya: keluarga
gandanya $\bigl(xy\,\P(X = x)\P(Y = y)\bigr)_{(x,y)}$ \omterm{def:b2:series:summable}{terjumlahkan}, sebab menurut
Fubini bagi keluarga (\cref{ch:b2:series})
\[
\sum_{x, y}\abs x \abs y\,\P(X{=}x)\P(Y{=}y)
= \Bigl(\sum_x \abs x \P(X{=}x)\Bigr)
\Bigl(\sum_y \abs y \P(Y{=}y)\Bigr) < \infty ;
\]
lalu berkat kebebasannya keluarga ini persis $\bigl(xy\,\P(X = x, Y =
y)\bigr)$, yang jumlahnya adalah $\E(XY)$ menurut pemindahan yang diterapkan pada
peubah $(X, Y) \mapsto xy$; dan Fubini lagi menilai jumlah tak
bertandanya sebagai hasil kali $\E(X)\E(Y)$.
\end{proof}

\begin{example}[Hasil kali, dengan dan tanpa
kebebasan]\label{ex:b2:randomvar:productcontrast}
Lemparkan dua dadu setimbang. Bila $Y$ adalah dadu kedua (yang bebas dari
yang pertama), maka $\E(XY) = \E(X)\E(Y) = 3.5^2 = 12.25$. Sedangkan bila
$Y = X$ (yakni ``hasil kali'' sebuah dadu dengan dirinya sendiri),
\[
\E(X^2) = \frac{1 + 4 + 9 + 16 + 25 + 36}{6} = \frac{91}{6}
\approx 15.17 \neq 12.25 :
\]
jadi \omterm{def:b2:randomvar:law}{hukum} marginal yang sama pada kedua skenarionya, \omterm{def:b2:randomvar:law}{hukum} bersama yang berbeda,
dan \omterm{def:b2:randomvar:expectation}{nilai harapan} hasil kali yang berbeda. Adapun pelajarannya, yang layak dipahat:
$\E(XY)$ adalah fungsional \emph{pasangannya}, bukan kedua
marginalnya --- dan senjangnya $\E(X^2) - \E(X)^2 \approx 2.92$
persis, menurut König--Huygens, merupakan \omterm{def:b2:randomvar:variance}{ragam}
$\frac{35}{12}$ dadunya.
\end{example}

\section{Ragam, kovariansi, dan ketaksamaan klasiknya}

\begin{definition}[Momen, ragam]\label{def:b2:randomvar:variance}
Peubah $X$ disebut punya \emph{momen orde 2} bila $X^2$ punya \omterm{def:b2:randomvar:expectation}{nilai harapan}
(dan lalu $X$ juga, berkat dominasi: sebab $\abs X \leq \frac{1 +
X^2}{2}$). Adapun \emph{ragam}\index{ragam} dan \emph{simpangan
baku}\index{simpangan baku} peubahnya lalu
\[
V(X) = \E\bigl((X - \E(X))^2\bigr)
= \E(X^2) - \E(X)^2 ,
\qquad
\sigma(X) = \sqrt{V(X)} ,
\]
(yakni bentuk keduanya --- rumus \emph{König--Huygens} --- lewat
menguraikan kuadratnya lalu memakai kelinearannya:
\[
\E\bigl((X - \E X)^2\bigr)
= \E\bigl(X^2 - 2X\,\E X + \E(X)^2\bigr)
= \E(X^2) - 2\,\E(X)^2 + \E(X)^2 ,
\]
dengan suku tengahnya memakai bahwa $\E X$ sebuah konstanta). Adapun untuk $X, Y$ yang bermomen kedua,
\emph{kovariansi}\index{kovariansi} keduanya adalah
\[
\operatorname{Cov}(X, Y)
= \E\bigl((X - \E X)(Y - \E Y)\bigr)
= \E(XY) - \E(X)\E(Y) .
\]
\end{definition}

\begin{theorem}[Perkakas ragam]\label{thm:b2:randomvar:variancerules}
Untuk peubah yang bermomen kedua:
\begin{enumerate}
\item $V(aX + b) = a^2\,V(X)$;
\item $V(X + Y) = V(X) + V(Y) + 2\operatorname{Cov}(X, Y)$, dan
lebih umum
\[
V\Bigl(\sum_{i=1}^n X_i\Bigr)
= \sum_{i=1}^n V(X_i)
+ 2\sum_{i < j}\operatorname{Cov}(X_i, X_j) ;
\]
\item bila $X, Y$ \omterm{def:b2:proba:independence}{saling bebas}, maka $\operatorname{Cov}(X, Y) = 0$
(sedangkan konversnya salah), jadi \omterm{def:b2:randomvar:variance}{ragam} peubah yang \omterm{def:b2:proba:independence}{saling bebas}
bertambah.
\end{enumerate}
\end{theorem}

\begin{proof}
Butir \emph{1} dan \emph{2} merupakan penguraian kuadrat ditambah kelinearannya;
sedangkan hasil kali $X_iX_j$ punya \omterm{def:b2:randomvar:expectation}{nilai harapan} berkat Cauchy--Schwarz di bawah
(atau berkat $\abs{X_iX_j} \leq \frac{X_i^2 + X_j^2}{2}$). Adapun butir \emph{3} adalah
\cref{thm:b2:randomvar:product} yang diterapkan pada peubah yang
dipusatkan. Contoh tandingan baku bagi konversnya: $X$ yang berhukum seragam
pada $\{-1, 0, 1\}$ dan $Y = X^2$ tak berkorelasi
(sebab $\E(XY) = \E(X^3) = 0 = \E X \cdot \E Y$) tetapi jelas bergantungan.
\end{proof}

\begin{theorem}[Ketaksamaan Markov dan
Chebyshev]\label{thm:b2:randomvar:markov}
\leavevmode
\begin{enumerate}
\item (Markov) Bila $X \geq 0$ punya \omterm{def:b2:randomvar:expectation}{nilai harapan}, maka untuk setiap $a
> 0$:
\[
\P(X \geq a) \leq \frac{\E(X)}{a} .
\]
\item (Chebyshev) Bila $X$ bermomen kedua, maka untuk setiap
$\varepsilon > 0$:
\[
\P\bigl(\abs{X - \E(X)} \geq \varepsilon\bigr)
\leq \frac{V(X)}{\varepsilon^2} .
\]
\end{enumerate}
\end{theorem}

\begin{proof}
\emph{1.} \omterm{def:b2:funcseq:def}{Titik demi titik} berlaku $a\,\mathbf{1}_{X \geq a} \leq X$ (sebab pada
\omterm{def:b2:proba:space}{kejadiannya} ruas kirinya $a \leq X$; sedangkan di luarnya $0 \leq X$). Lalu ambillah
\omterm{def:b2:randomvar:expectation}{nilai harapannya}: $a\,\P(X \geq a) \leq \E(X)$ berkat kemonotonannya dan
$\E(\mathbf{1}_A) = \P(A)$.
\emph{2.} Terapkan Markov pada peubah taknegatif $(X - \E
X)^2$ di aras $a = \varepsilon^2$: sebab \omterm{def:b2:proba:space}{kejadian} $\{(X - \E X)^2
\geq \varepsilon^2\}$ persis merupakan $\{\abs{X - \E X} \geq
\varepsilon\}$.
\end{proof}

\begin{example}[Tak berkorelasi tetapi terlem
bersama]\label{ex:b2:randomvar:sumdiff}
Lemparkan dua dadu setimbang, dengan $X$ dan $Y$ \omterm{def:b2:proba:independence}{saling bebas}, lalu tetapkan $S = X +
Y$ dan $D = X - Y$. Menurut kebilinearan \omterm{def:b2:randomvar:variance}{kovariansinya},
\[
\operatorname{Cov}(S, D) = V(X) - V(Y) +
\operatorname{Cov}(Y, X) - \operatorname{Cov}(X, Y) = V(X) -
V(Y) = 0 :
\]
jadi jumlah dan selisihnya tak berkorelasi. \omterm{def:b2:proba:independence}{Saling bebas}? Tentu
tidak: sebab $S = 12$ memaksa $D = 0$, sedangkan $\P(D = 0) =
\frac16$ tanpa syarat. Jadi korelasi hanya menguji bagian
\emph{linear} sebuah kebergantungan; sebab di sini kebergantungannya
diusung kendala bahwa $S$ dan $D$ berparitas sama,
yang tak terlihat oleh \omterm{def:b2:randomvar:variance}{kovariansinya}. (Adapun bagi pasangan ini, \omterm{def:b2:randomvar:variance}{kovariansi}
nolnya menuntut $V(X) = V(Y)$: jadi \omterm{def:b2:randomvar:law}{distribusi} yang identik,
bukan kebebasan, yang mengerjakannya.)
\end{example}

\begin{example}[Ketika Markov eksak]\label{ex:b2:randomvar:markovsharp}
Ketaksamaan Markov menjadi kesamaan persis ketika tak ada yang
terbuang pada batas $a\,\mathbf 1_{X\geq a} \leq X$: yakni
peubahnya hanya boleh mengambil nilai $0$ dan $a$. Secara konkret,
bila $\P(X = a) = \pi$ dan $\P(X = 0) = 1 - \pi$, maka $\E(X)
= a\pi$ dan
\[
\P(X \geq a) = \pi = \frac{\E(X)}{a} .
\]
Adapun pembacaan realistisnya: pada sebuah populasi yang kekayaan rata-ratanya
$100$ dan kekayaannya entah $0$ entah $10^6$, proporsi
jutawannya persis $10^{-4}$ --- yakni batas Markov, yang tercapai
persis lewat ketaksamaan maksimalnya. Jadi setiap kali $X$ menyebar atas
nilai antaranya batasnya menjadi tegas, kerap secara liar;
tetapi seperti ditunjukkan kasus ekstremnya, tak ada ketaksamaan yang lebih baik yang dapat
disarikan dari rerataannya sendirian.
\end{example}

\begin{example}[Chebyshev bersifat tajam --- tanpa hipotesis
lebih lanjut]\label{ex:b2:randomvar:chebsharp}
Tetapkan $\varepsilon > 0$ dan $q \in \intoc01$, lalu biarkanlah $X$ mengambil
nilai $\pm\varepsilon$ masing-masing berpeluang $\frac q2$ dan
$0$ berpeluang $1 - q$. Maka $\E(X) = 0$, $V(X) =
q\varepsilon^2$, dan
\[
\P\bigl(\abs{X - \E X} \geq \varepsilon\bigr) = q
= \frac{V(X)}{\varepsilon^2} :
\]
yakni kesamaan pada Chebyshev. Jadi ketaksamaannya tak dapat diperbaiki
dengan memakai \omterm{def:b2:randomvar:variance}{ragamnya} saja --- sebab peluruhan $1/\varepsilon^2$-nya adalah
harga persis informasi momen keduanya. Adapun peluruhan yang lebih cepat
menuntut hipotesis yang lebih kuat: sebab keterbatasan peubahnya
membeli pemusatan \emph{eksponensial}, sebagaimana dipratinjau
\cref{exo:b2:randomvar:7} dan dikembangkan secara sistematis
oleh soal akhir pekan bab ini.
\end{example}

\begin{theorem}[Cauchy--Schwarz dan Jensen]\label{thm:b2:randomvar:jensen}
\leavevmode
\begin{enumerate}
\item (Cauchy--Schwarz) Bila $X, Y$ bermomen kedua, maka $XY$ punya
\omterm{def:b2:randomvar:expectation}{nilai harapan} dan $\E(XY)^2 \leq \E(X^2)\,\E(Y^2)$; sehingga
$\operatorname{Cov}(X,Y)^2 \leq V(X)V(Y)$.
\item (Jensen) Bila $\varphi \colon I \to \R$ \omterm{def:b2:affine:convex}{cembung} pada sebuah
interval yang memuat $X(\Omega)$, dan $X$ serta $\varphi(X)$ punya
\omterm{def:b2:randomvar:expectation}{nilai harapan}, maka
\[
\varphi\bigl(\E(X)\bigr) \leq \E\bigl(\varphi(X)\bigr) .
\]
\end{enumerate}
\end{theorem}

\begin{proof}
\emph{1.} Keterjumlahan $XY$: sebab $\abs{XY} \leq \frac{X^2 + Y^2}2$.
Adapun pemetaan $(X, Y) \mapsto \E(XY)$ merupakan bentuk bilinear \omterm{def:b2:quadratic:adjoint}{simetrik} yang
positif pada ruang peubah yang bermomen kedua, jadi ketaksamaan
Cauchy--Schwarz abstrak pada \cref{ch:b2:quadratic} berlaku
(sebab \emph{semi}definit positif sudah cukup bagi ketaksamaannya).
Lalu menerapkannya pada peubah yang dipusatkan memberi batas \omterm{def:b2:randomvar:variance}{kovariansinya}.

\emph{2.} Pertama, $m = \E(X)$ terletak di $I$: sebab $I$ merupakan interval
yang memuat semua nilai $X$, dan \omterm{def:b2:randomvar:expectation}{nilai harapan} bersifat monoton, jadi $m$
berada di antara $\inf X(\Omega)$ dan $\sup X(\Omega)$. Lalu menurut
teorema garis penyangga bagi fungsi \omterm{def:b2:affine:convex}{cembung}
(\cref{ch:b2:realfun}), ada $\alpha, \beta$ dengan
$\varphi(t) \geq \alpha t + \beta$ untuk setiap $t \in I$ dan
$\varphi(m) = \alpha m + \beta$. Maka, \omterm{def:b2:funcseq:def}{titik demi titik} pada $\Omega$,
$\varphi(X) \geq \alpha X + \beta$; jadi setelah \omterm{def:b2:randomvar:expectation}{nilai harapannya} diambil,
\[
\E\bigl(\varphi(X)\bigr) \geq \alpha\,\E(X) + \beta
= \varphi\bigl(\E(X)\bigr). \qedhere
\]
\end{proof}

\begin{example}
Jensen dengan $\varphi(t) = t^2$ memberi $\E(X)^2 \leq \E(X^2)$ ---
yakni kepositifan \omterm{def:b2:randomvar:variance}{ragamnya}; sedangkan dengan $\varphi(t) = 1/t$ pada
$\intoo{0}{\infty}$: $\frac{1}{\E X} \leq \E\bigl(\frac1X\bigr)$
--- yakni rerata harmoniknya di bawah rerata aritmetikanya, kini dalam bentuk
acak.
\end{example}

\begin{remark}[Jebakan yang sering muncul]
(i) Rumus $\E(XY) = \E(X)\E(Y)$ \emph{menuntut} kebebasan (atau sekurangnya
\omterm{def:b2:randomvar:variance}{kovariansi} nol): sebab mengambil $Y = X$ memberi $\E(X^2) \neq
\E(X)^2$ setiap kali $V(X) > 0$. (ii) Demikian pula $V(X + X) =
4V(X)$, bukan $2V(X)$: sebab \omterm{def:b2:randomvar:variance}{ragam} hanya bertambah antar suku yang \omterm{def:b2:proba:independence}{saling bebas}
(atau yang tak berkorelasi). (iii) Nilai $\E(f(X))$ bukanlah
$f(\E(X))$; dan untuk $f$ yang \omterm{def:b2:affine:convex}{cembung} Jensen bahkan memberitahukanmu
arah galatnya, seperti pada
\cref{ex:b2:randomvar:transferex}. (iv) Keberadaannya sebuah
hipotesis yang sejati: sebab untuk peubah St Petersburg $X = 2^K$ dengan
$\P(K = k) = 2^{-k}$ (untuk $k \geq 1$),
\[
\sum_{k\geq1}2^k\cdot2^{-k} = \sum_{k\geq1}1 = \infty :
\]
jadi $X$ berhingga hampir pasti namun tak punya \omterm{def:b2:randomvar:expectation}{nilai harapan}, sehingga tak ada
harga masuk yang adil bagi permainannya. Jadi keterjumlahan pada
definisi $\E$ bukanlah kerewelan tata buku --- melainkan tempat
ekor yang berat terdeteksi. (v) Akhirnya, teorema
pemindahannya menuntut keterjumlahan \emph{\omterm{def:b2:series:def}{mutlak}} sebelum sembarang
penataan ulang jumlahnya atas nilainya menjadi sah
(\cref{ch:b2:series}).
\end{remark}


\begin{example}[Chebyshev pada seratus
lemparan]\label{ex:b2:randomvar:hundredtosses}
Untuk $X \sim \mathcal B(100, \frac12)$: berlaku $\E X = 50$ dan $V(X) =
25$. Lalu Chebyshev dengan $\varepsilon = 6$:
\[
\P(45 \leq X \leq 55) = \P(\abs{X - 50} < 6)
\geq 1 - \frac{25}{36} \approx 0.31 ,
\]
sedangkan jumlah binomial \omterm{def:b2:multint:exact}{eksaknya} memberi $\approx 0.73$. Jadi
jaminan $31\%$-nya jauh dari kebenarannya, tetapi ia hanya menuntut
rerataannya dan \omterm{def:b2:randomvar:variance}{ragamnya} \emph{saja} --- sebab sertifikat yang sama
berlaku kata demi kata bagi sembarang peubah dengan $\E = 50$ dan $V = 25$,
seeksotis apa pun, dan \cref{ex:b2:randomvar:chebsharp} menunjukkan
bahwa suatu peubah semacam itu menjenuhkannya. Jadi keuniversalannya berharga;
sedangkan ketika \omterm{def:b2:randomvar:law}{distribusinya} sungguh binomial, perkakas eksponensial
pada soal akhir pekan menutup sebagian besar senjangnya.
\end{example}

\begin{example}[Korelasi sebuah bagian dengan
keseluruhannya]\label{ex:b2:randomvar:partwhole}
Untuk $X, Y$ yang \omterm{def:b2:proba:independence}{saling bebas} dan berdistribusi identik dengan \omterm{def:b2:randomvar:variance}{ragam}
$\sigma^2 > 0$, seberapa berkorelasikah satu sukunya dengan jumlahnya
$S = X + Y$? Hitunglah
\[
\operatorname{Cov}(X, S) = \operatorname{Cov}(X, X) +
\operatorname{Cov}(X, Y) = \sigma^2 + 0 = \sigma^2,
\qquad V(S) = 2\sigma^2,
\]
jadi koefisien korelasinya adalah
\[
\rho(X, S) = \frac{\operatorname{Cov}(X,
S)}{\sigma(X)\,\sigma(S)}
= \frac{\sigma^2}{\sigma\cdot\sigma\sqrt2}
= \frac{1}{\sqrt2} \approx 0.707 ,
\]
apa pun \omterm{def:b2:randomvar:law}{hukum} bersamanya --- dadu, koin, atau cacah Poisson. Adapun dengan
$n$ suku perhitungan yang sama memberi $\rho(X_1, S_n) =
1/\sqrt n$: jadi pengaruh tiap suku individualnya pada totalnya
mengencer seperti akar kuadrat, dan itulah bayangan korelasional
bagi skala fluktuasi $\sqrt n$. Adapun Cauchy--Schwarz
menjamin $\abs\rho \leq 1$ selalu; dan di sini batasnya tercapai
persis pada kasus merosot $n = 1$ lalu meluruh secara terduga
sesudahnya.
\end{example}

\begin{example}[AM--GM terbobot dari
Jensen]\label{ex:b2:randomvar:amgm}
Misalkan $Y$ mengambil nilai positif $a_1, \dots, a_k$ dengan
peluang $\lambda_1, \dots, \lambda_k$. Adapun fungsi
$-\ln$ \omterm{def:b2:affine:convex}{cembung} pada $\intoo0\infty$, jadi Jensen memberi
$-\ln\E(Y) \leq \E(-\ln Y)$, yakni
\[
a_1^{\lambda_1}a_2^{\lambda_2}\cdots a_k^{\lambda_k}
\;\leq\; \lambda_1a_1 + \lambda_2a_2 + \dots + \lambda_ka_k :
\]
yaitu ketaksamaan aritmetika--geometri yang terbobot, dengan kesamaan
bila dan hanya bila $Y$ tetap. Adapun bobot yang sama $\lambda_i = \frac1k$
memulihkan AM--GM klasiknya. Jadi peluang diam-diam telah membuktikan
sebuah teorema yang murni aljabar: sebab memilih sebuah \omterm{def:b2:randomvar:law}{hukum} peluang
tak lain alat tata buku bagi kombinasi \omterm{def:b2:affine:convex}{cembung} --- yakni
sudut pandang barisentrik \cref{ch:b2:affine} sekali lagi, kini
dengan Jensen sebagai mesinnya.
\end{example}

\section{Hukum bilangan besar yang lemah}

\begin{theorem}[Hukum bilangan besar yang lemah]\label{thm:b2:randomvar:lln}
Misalkan $(X_k)_{k \geq 1}$ \omterm{def:b2:randomvar:law}{peubah acak} yang \omterm{def:b2:proba:independence}{saling bebas} berpasangan
dan berhukum sama, serta bermomen kedua; lalu tulislah $m = \E(X_1)$
dan $S_n = X_1 + \dots + X_n$. Maka untuk setiap $\varepsilon > 0$:
\[
\P\Bigl(\,\Bigl|\frac{S_n}{n} - m\Bigr| \geq \varepsilon\Bigr)
\;\leq\; \frac{V(X_1)}{n\,\varepsilon^2}
\xrightarrow[n \to \infty]{} 0 .
\]
\end{theorem}

\begin{proof}
Berkat kelinearannya $\E(S_n/n) = m$; lalu menurut
\cref{thm:b2:randomvar:variancerules} (sebab kebebasan berpasangannya
membunuh \omterm{def:b2:randomvar:variance}{kovariansinya}) $V(S_n) = n\,V(X_1)$, jadi $V(S_n/n) =
V(X_1)/n$. Sehingga ketaksamaan Chebyshev yang diterapkan pada $S_n/n$ memberi
batasnya.
\end{proof}

\begin{remark}
Inilah teorema yang menghubungkan peluang dengan frekuensi: sebab untuk
$X_k$ berupa indikator sebuah \omterm{def:b2:proba:space}{kejadian} $A$ pada pengulangan yang \omterm{def:b2:proba:independence}{saling bebas},
$S_n/n$ merupakan frekuensi $A$ yang teramati, dan \omterm{def:b2:randomvar:law}{hukum} bilangan
besar mengatakan bahwa ia memusat di sekitar $\P(A)$ dengan laju
$\frac{p(1-p)}{n\varepsilon^2}$. Adapun \omterm{def:b2:randomvar:law}{hukum} yang \emph{kuat} (yakni $S_n/n \to
m$ hampir pasti) merupakan teorema Tahun ke-3 --- namun buktinya untuk
momen keempat sudah terjangkau: lihat \cref{exo:b2:randomvar:9},
yang menjalankan Borel--Cantelli pada batas bertipe Chebyshev. Adapun
taksiran Chebyshev yang sama menggerakkan bukti \omterm{thm:b2:funcseq:weierstrass}{polinomial Bernstein} bagi
teorema hampiran Weierstrass pada \cref{ch:b2:funcseq} --- sebab
lema pencacahan di sana \emph{adalah} \omterm{def:b2:randomvar:law}{hukum} bilangan besar yang lemah
yang menyamar.
\end{remark}

\begin{example}[Mengumpulkan lima puluh kupon]
\label{ex:b2:randomvar:couponnumeric}
Pengumpul kupon pada \cref{exo:b2:randomvar:3} dengan $n =
50$ mainan berbeda: maka total yang diharapkan adalah
\[
\E(T_{50}) = 50\,H_{50} = 50\sum_{k=1}^{50}\frac1k
\approx 50 \times 4.499 \approx 225
\]
kotak --- yakni empat setengah kali tebakan naifnya $50$. Adapun
pertumbuhan harmoniknya seluruh ceritanya: sebab $25$ mainan pertamanya
tiba dalam sekitar $50\ln2 \approx 35$ kotak, sedangkan mainan
\emph{terakhirnya} sendirian berharga $50$ kotak secara rata-rata (yakni
tunggu geometri berparameter $\frac1{50}$). Jadi masalah
penyelesaian didominasi babak akhirnya, dan itulah sebabnya
\cref{exo:b2:randomvar:12} menemukan fluktuasi berorde $n$
--- yakni besarnya tunggu geometri terakhir itu --- di sekitar
rerataannya $n\ln n$.
\end{example}

\begin{example}[Seberapa besar $n$ harus?]\label{ex:b2:randomvar:llnnumeric}
Untuk memaku frekuensi teramatinya dalam $\varepsilon = 0.01$ dari
$\P(A)$ dengan kepercayaan $95\%$, batas Chebyshev menuntut
\[
\frac{p(1-p)}{n\varepsilon^2} \leq \frac{1}{4n\varepsilon^2}
\leq 0.05,
\qquad\text{yakni}\qquad
n \geq \frac{1}{4\cdot0.05\cdot(0.01)^2} = 50\,000 .
\]
Kebergantungannya brutal dalam $\varepsilon$ (yakni kuadratik) dan
lunak dalam kepercayaannya (yakni linear dalam $1/\alpha$). Adapun kedua ciri
itu merupakan sifat \emph{batasnya}, bukan sifat kebenarannya: sebab
ketaksamaan eksponensial pada soal akhir pekan menurunkan harga
kepercayaannya dari $1/\alpha$ menjadi $\ln(1/\alpha)$ --- jadi spesifikasi
yang sama kelak berharga sekitar $18\,500$ cuplikan di sana ---
sedangkan skala $1/\varepsilon^2$-nya sejati dan
tak terperbaiki. Jadi mengetahui bagian mana sebuah batas yang longgar sama
bergunanya dengan batasnya sendiri.
\end{example}

\begin{omfigure}
\begin{tikzpicture}[x=5.5cm, y=1.05cm]
  \draw[->] (-0.05,0) -- (1.1,0) node[below] {$S_n/n$};
  \draw[->] (0,-0.1) -- (0,2.35);
  \draw[dashed] (0.4,0) -- (0.4,2.15);
  \draw[dashed] (0.6,0) -- (0.6,2.15);
  \node[below, font=\scriptsize] at (0.4,0)
    {$m{-}\varepsilon$};
  \node[below, font=\scriptsize] at (0.5,0) {$m$};
  \node[below, font=\scriptsize] at (0.6,0)
    {$m{+}\varepsilon$};
  \draw[omDef, thick] plot[domain=0.06:0.94, samples=80]
    (\x, {0.8*exp(-(\x-0.5)^2/0.0288)});
  \draw[omThm, very thick] plot[domain=0.32:0.68, samples=80]
    (\x, {2.0*exp(-(\x-0.5)^2/0.0032)});
  \node[omDef, font=\small] at (0.16, 0.85) {$n$ kecil};
  \node[omThm, font=\small] at (0.71, 1.95) {$n$ besar};
\end{tikzpicture}

{\small \omterm{def:b2:randomvar:law}{Hukum} bilangan besar sebagai gambar: sebab \omterm{def:b2:randomvar:law}{hukum}
$S_n/n$ (yang digambar secara skematis) mempertahankan pusatnya $m$ tetapi
menyempit ketika $n$ tumbuh, jadi peluang di luar pita
$\intcc{m-\varepsilon}{m+\varepsilon}$ --- yakni kedua ekornya ---
menciut ke nol. Adapun Chebyshev membatasi ekornya oleh
$V(X_1)/(n\varepsilon^2)$; sedangkan soal akhir pekan menunjukkan bahwa keduanya
sebenarnya kecil secara eksponensial.}
\end{omfigure}

\begin{remark}[Pandangan ke depan di dalam jilid ini]
Ke depan, segalanya di sini menyuapi \cref{ch:b2:genfun}: sebab
\omterm{def:b2:randomvar:expectation}{nilai harapan} $\E(t^X)$ bagi satu fungsi $X$ yang cerdik memampatkan
seluruh \omterm{def:b2:randomvar:law}{hukumnya} menjadi sebuah deret pangkat, momennya menjadi
turunan di $1$, dan kesamaan bertipe Wald bagi jumlah acak
mengusung teori proses bercabangnya; sedangkan teorema hasil kali bagi
peubah yang \omterm{def:b2:proba:independence}{saling bebas} menjadi kemultiplikatifan \omterm{ex:b2:powerseries:fibonacci}{fungsi pembangkitnya}. Ke belakang,
\omterm{def:b2:randomvar:expectation}{nilai harapan} adalah \omterm{def:b2:affine:barycenter}{barisentrum} dengan bobot peluang (\cref{ch:b2:affine}),
ketaksamaan Jensen adalah geometri garis penyangga bagi fungsi \omterm{def:b2:affine:convex}{cembung}
(\cref{ch:b2:realfun}), dan metode momen eksponensial pada
soal akhir pekan bab ini adalah Markov yang diterapkan pada
$\eu^{tX}$ --- jadi satu ketaksamaan, yang ditingkatkan oleh satu penggantian
peubah yang baik, membentang tiga bab.
\end{remark}

\section{Latihan}

\begin{exercise}[$\star$]\label{exo:b2:randomvar:1}
Hitunglah $\E(X)$ dan $V(X)$ untuk $X \sim \mathcal{B}(n, p)$ (lewat
indikator), $X \sim \mathcal{P}(\lambda)$ (tunjukkan $V(X) =
\lambda$), dan $X \sim \mathcal{G}(p)$ (tunjukkan $V(X) =
\frac{1-p}{p^2}$; pakailah $\E(X(X-1))$ beserta turunan kedua
deret geometrinya).
\end{exercise}

\begin{exercise}[$\star$]\label{exo:b2:randomvar:2}
Misalkan $X \sim \mathcal{P}(\lambda)$ dan $Y \sim \mathcal{P}(\mu)$
\omterm{def:b2:proba:independence}{saling bebas}. Tunjukkan bahwa $X + Y \sim \mathcal{P}(\lambda + \mu)$
(lewat konvolusi bobotnya dan teorema binomial), dan bahwa
\omterm{def:b2:randomvar:law}{hukum} bersyarat $X$ jika diketahui $X + Y = n$ bersifat binomial
$\mathcal{B}\bigl(n, \frac{\lambda}{\lambda + \mu}\bigr)$.
\end{exercise}

\begin{exercise}[$\star$]\label{exo:b2:randomvar:3}
(Pengumpul kupon, \omterm{def:b2:randomvar:expectation}{nilai harapannya}) Sebuah merek sereal menyembunyikan satu dari $n$
mainan berbeda, secara seragam, di tiap kotaknya. Misalkan $T_n$ banyaknya
kotak yang diperlukan untuk mengumpulkan seluruh $n$ mainannya. Dengan menulis $T_n$ sebagai jumlah
peubah geometri yang \omterm{def:b2:proba:independence}{saling bebas} (yakni waktu melihat mainan \emph{baru} ketika
$k$ masih hilang), tunjukkanlah
\[
\E(T_n) = n\sum_{k=1}^{n}\frac{1}{k} \sim n\ln n
\]
(yang setara lewat pembandingan deret--integral pada
\cref{ch:b2:comparison}).
\end{exercise}

\begin{exercise}[$\star\star$]\label{exo:b2:randomvar:4}
Misalkan $X \geq 0$ bernilai bulat. Buktikan \emph{rumus ekornya}
\[
\E(X) = \sum_{n=1}^{\infty} \P(X \geq n)
\]
(ketika salah satu ruasnya berhingga), dengan menulis $X =
\sum_{n\geq1}\mathbf{1}_{X \geq n}$ lalu menukarkan penjumlahannya
(lewat Fubini bagi keluarga taknegatif). Lalu perolehlah $\E(X) = \frac1p$ bagi
\omterm{def:b2:randomvar:law}{hukum} geometrinya.
\end{exercise}

\begin{exercise}[$\star\star$]\label{exo:b2:randomvar:5}
(Pencuplikan tanpa pengembalian lebih memusat) Sebuah guci berisi $N$
bola, dan $M$ di antaranya putih. Tariklah $n \leq N$ tanpa pengembalian lalu
misalkan $X$ mencacah yang putihnya (yakni \omterm{def:b2:randomvar:law}{hukum} \emph{hipergeometrik}). Dengan memakai
indikator $X = \sum_{i=1}^n Y_i$ dengan $Y_i$ menyatakan tarikan ke-$i$:
tunjukkanlah bahwa tiap $Y_i$ bersifat Bernoulli berparameter $p = M/N$ (berkat kesimetriannya!),
lalu simpulkan $\E(X) = np$ persis seperti dengan pengembalian, kemudian tunjukkan
$\operatorname{Cov}(Y_i, Y_j) = -\frac{p(1-p)}{N-1} < 0$ untuk $i
\neq j$, sehingga $V(X) = np(1-p)\frac{N - n}{N - 1} \leq np(1-p)$.
\end{exercise}

\begin{exercise}[$\star\star$]\label{exo:b2:randomvar:6}
Misalkan $X$ bermomen kedua. Tunjukkan bahwa $c \mapsto \E\bigl((X -
c)^2\bigr)$ minimum persis di $c = \E(X)$, dengan minimum
$V(X)$. Lalu tunjukkan bahwa $\P(X = \E(X)) = 1$ bila dan hanya bila $V(X) =
0$. \emph{(Untuk butir keduanya: bila $V(X) = 0$, pakailah Chebyshev
dengan $\varepsilon = 1/n$ beserta kekontinuan monoton,
\cref{thm:b2:proba:continuity}.)}
\end{exercise}

\begin{exercise}[$\star\star\star$]\label{exo:b2:randomvar:7}
(Pemusatan mengalahkan Markov) Misalkan $S_n \sim \mathcal{B}(n,
\frac12)$ (yakni banyaknya gambar pada $n$ lemparan setimbang). Bandingkanlah
batas yang diberikan Markov (yakni $\P(S_n \geq \frac{3n}{4})$), Chebyshev,
dan metode eksponensial (Chernoff):
\[
\P\Bigl(S_n \geq \frac{3n}4\Bigr)
\leq \E\bigl(e^{tS_n}\bigr)e^{-3nt/4}
= \Bigl(\frac{1 + e^t}{2}\Bigr)^n e^{-3nt/4}
\quad (t > 0),
\]
lalu optimumkan $t$ untuk memperoleh batas yang kecil secara eksponensial. \emph{(Di
$t = \ln 3$: batasnya $\bigl(2\cdot 3^{-3/4}\bigr)^n \approx
(0.877)^n$.)}
\end{exercise}

\begin{exercise}[$\star\star\star$]\label{exo:b2:randomvar:8}
(Weierstrass lagi, secara peluang) Misalkan $f \colon [0,1] \to \R$
\omterm{def:b2:metric:continuity}{kontinu} dan $S_n \sim \mathcal{B}(n, x)$. Tunjukkan bahwa
\omterm{thm:b2:funcseq:weierstrass}{polinomial Bernstein} $B_nf(x) = \sum_{k=0}^n f\bigl(\frac
kn\bigr)\binom nk x^k(1-x)^{n-k}$ sama dengan
$\E\bigl[f\bigl(\frac{S_n}{n}\bigr)\bigr]$, lalu turunkan ulang
taksiran $\abs{B_nf(x) - f(x)} \leq \omega_f(\delta) +
\frac{2\norm f_\infty}{4n\delta^2}$ pada \cref{ch:b2:funcseq} dalam
bahasa peluang ini (yakni pecahlah pada $\bigl|\frac{S_n}{n} -
x\bigr| \geq \delta$ lalu pakailah Chebyshev).
\end{exercise}

\begin{exercise}[$\star\star\star$]\label{exo:b2:randomvar:9}
(\omterm{def:b2:randomvar:law}{Hukum} kuat di bawah momen keempat) Misalkan $(X_k)$ \omterm{def:b2:proba:independence}{saling bebas},
berdistribusi identik, terpusat (dengan $\E X_1 = 0$), dan
$\E(X_1^4) < \infty$. Dengan menguraikan $\E(S_n^4)$ lalu mencacah
suku yang bertahan (yakni hanya suku $\E(X_i^4)$ dan $\E(X_i^2X_j^2)$ dengan $i
\neq j$), tunjukkanlah $\E(S_n^4) \leq C n^2$ untuk sebuah konstanta $C$. Lalu turunkan
$\sum_n \P\bigl(\abs{S_n/n} \geq \varepsilon\bigr) < \infty$
untuk tiap $\varepsilon > 0$ (lewat Markov pada orde 4) kemudian rampungkan dengan
Borel--Cantelli (\cref{thm:b2:proba:borelcantelli}) bahwa
$S_n/n \to 0$ hampir pasti dalam perumusan yang sesuai: yakni bahwa
\omterm{def:b2:proba:space}{kejadian} $\bigcap_{j}\bigcup_N\bigcap_{n \geq N}\{\abs{S_n/n} <
\frac1j\}$ berpeluang $1$.
\end{exercise}

\begin{exercise}[$\star$]\label{exo:b2:randomvar:10}
Dua dadu setimbang dilempar; misalkan $M$ yang lebih besar di antara kedua
hasilnya. Dengan memakai rumus ekor pada \cref{exo:b2:randomvar:4}
(versi berhingganya), tunjukkanlah
\[
\E(M) = \sum_{k=1}^{6}\P(M \geq k)
= 6 - \sum_{j=0}^5\Bigl(\frac j6\Bigr)^2 = \frac{161}{36}
\approx 4.47 .
\]
\end{exercise}

\begin{exercise}[$\star\star$]\label{exo:b2:randomvar:11}
Misalkan $F_n$ banyaknya titik tetap sebuah permutasi acak seragam
atas $\{1, \dots, n\}$ (dengan $n \geq 2$). Dengan menulis $F_n =
\sum_i\mathbf 1_{\sigma(i) = i}$, hitunglah $\E(F_n) = 1$ dan
$\operatorname{Cov}(\mathbf 1_{\sigma(i)=i}, \mathbf
1_{\sigma(j)=j}) = \frac1{n^2(n-1)}$ untuk $i \neq j$, lalu
simpulkan $V(F_n) = 1$: jadi rata-rata satu huruf tetap, dengan
\omterm{def:b2:randomvar:variance}{ragam} persis $1$, berapa pun $n$.
\end{exercise}

\begin{exercise}[$\star\star\star$]\label{exo:b2:randomvar:12}
(Pengumpul kupon, pemusatannya) Pada latar
\cref{exo:b2:randomvar:3}, tunjukkanlah
\[
V(T_n) = \sum_{k=1}^n\frac{1 - k/n}{(k/n)^2}
\leq n^2\sum_{k=1}^n\frac{1}{k^2} \leq \frac{\pi^2}{6}n^2,
\]
dengan memakai kebebasan tahap geometrinya beserta $V(\mathcal
G(p)) = \frac{1-p}{p^2}$ (\cref{exo:b2:randomvar:1}; dan nilai
$\pi^2/6$-nya adalah \cref{ex:b2:fourier:basel}). Lalu turunkan dengan
Chebyshev bahwa $\dfrac{T_n}{n\ln n} \to 1$ \emph{dalam
peluang}: jadi total waktu pengumpulnya adalah $n\ln n$ hingga
fluktuasi berorde $n$.
\end{exercise}

\section{Soal: kotak perkakas pemusatan, dari Markov ke
Hoeffding}

\begin{problem}[{Soal akhir pekan --- pemusatan eksponensial
dengan tangan, dan berapa orang yang harus ditanya sebuah jajak pendapat}]
\label{pb:b2:randomvar:1}
Ketaksamaan Markov berharga satu momen dan membeli peluruhan $1/a$;
sedangkan Chebyshev berharga dua momen dan membeli $1/\varepsilon^2$ ---
dan \cref{ex:b2:randomvar:chebsharp} menunjukkan bahwa itulah semua yang dapat
dibeli momen tersebut. Soal ini mendaki sisa tangganya:
yakni metode eksponensial (Chernoff) dengan lajunya yang \emph{\omterm{def:b2:multint:exact}{eksak}}
bagi lemparan koin, ketaksamaan Hoeffding bagi semua peubah yang terbatas,
dan panennya --- yakni ukuran cuplikan yang gamblang dan jujur
bagi jajak pendapat, penetapan pemenang pemilu, dan pengujian koin.
\index{ketaksamaan pemusatan}\index{ketaksamaan
Hoeffding}\index{batas Chernoff} Di sepanjang soal ini, $S_n \sim
\mathcal B(n, p)$ merupakan jumlah $n$ peubah Bernoulli yang \omterm{def:b2:proba:independence}{saling
bebas} dan $\widehat p_n = S_n/n$ frekuensi empirisnya.

\medskip
\noindent\textbf{Bagian I --- Penakaran pada koin yang setimbang.}
Di sini $p = \frac12$ dan $a \in \intoo{\frac12}{1}$.
\begin{enumerate}
  \item Markov di aras $an$: tunjukkan $\P(S_n \geq an) \leq
        \frac1{2a}$, yakni batas yang bahkan tak menuju $0$.
        Lalu di mana Markov kehilangan sebanyak itu?
  \item Chebyshev: dengan memakai kesetangkupan binomial setimbangnya
        terhadap $n/2$, tunjukkanlah
        \[
        \P(S_n \geq an) = \tfrac12\,
        \P\bigl(\abs{S_n - \tfrac n2} \geq n(a -
        \tfrac12)\bigr)
        \leq \frac{1}{8n(a - 1/2)^2},
        \]
        yakni $\frac2n$ di $a = \frac34$: jadi peluruhan polinomial pada
        akhirnya.
  \item (Chernoff, aras umum) Hitunglah
        $\E(\eu^{tS_n}) = \bigl(\frac{1 + \eu^t}2\bigr)^n$ lalu
        optimumkan $\P(S_n \geq an) \leq
        \E(\eu^{tS_n})\eu^{-tan}$ atas $t > 0$: tunjukkan bahwa
        $t$ optimumnya adalah $\ln\frac{a}{1-a}$ dan
        \[
        \P(S_n \geq an) \leq \eu^{-n\,I(a)},
        \qquad
        I(a) = \ln 2 + a\ln a + (1-a)\ln(1-a) > 0 .
        \]
        Lalu periksalah bahwa $a = \frac34$ memulihkan batas
        $\bigl(2\cdot3^{-3/4}\bigr)^n$ pada
        \cref{exo:b2:randomvar:7}.
  \item (Eksponennya \omterm{def:b2:multint:exact}{eksak}) Misalkan $k = an$ bilangan bulat.
        Dari fakta bahwa $\binom nk a^k(1-a)^{n-k}$ merupakan
        yang terbesar di antara $n + 1$ suku sebuah \omterm{def:b2:randomvar:law}{distribusi}
        peluang, buktikanlah $\binom nk \geq
        \frac{\eu^{nH(a)}}{n+1}$ dengan $H(a) = -a\ln a -
        (1-a)\ln(1-a)$, lalu turunkan batas bawah yang bersesuaian
        \[
        \P(S_n \geq an) \geq \binom{n}{an}2^{-n}
        \geq \frac{\eu^{-n\,I(a)}}{n + 1} .
        \]
  \item Tabelkan ketiga batasnya di $n = 100$ dan $a =
        \frac34$: yakni Markov $\frac23$, Chebyshev $0.02$, dan
        Chernoff $\approx 2.1\cdot10^{-6}$ (sedangkan nilai benarnya
        $\approx 2.8\cdot10^{-7}$). Lalu pelajarannya, dalam satu kalimat?
\end{enumerate}

\medskip
\noindent\textbf{Bagian II --- Ketaksamaan Hoeffding.}
\begin{enumerate}[resume]
  \item (Kasus Rademacher) Untuk $\varepsilon = \pm1$ yang masing-masing
        berpeluang $\frac12$, buktikanlah
        \[
        \E(\eu^{t\varepsilon}) = \cosh t \leq \eu^{t^2/2}
        \qquad (t \in \R)
        \]
        dengan membandingkan kedua deretnya suku demi suku (sebab $(2k)! \geq
        2^kk!$).
  \item Turunkan, untuk peubah Rademacher yang \omterm{def:b2:proba:independence}{saling bebas}
        $\varepsilon_1, \dots, \varepsilon_n$ dan setiap $s >
        0$:
        \[
        \P\Bigl(\sum_{i=1}^n\varepsilon_i \geq s\Bigr)
        \leq \eu^{-s^2/(2n)} .
        \]
  \item Terjemahkan ke koin setimbang (dengan $X_i =
        \frac{1+\varepsilon_i}2$): $\P\bigl(\widehat p_n -
        \tfrac12 \geq \delta\bigr) \leq \eu^{-2n\delta^2}$,
        beserta versi dua sisinya dengan faktor $2$.
  \item (Lema Hoeffding) Misalkan $X \in \intcc01$ dengan $\E X =
        p$, dan $\psi(t) = \ln\E(\eu^{tX})$. Benarkanlah bahwa
        $\psi$ \omterm{def:b2:diffcalc:differential}{terdiferensialkan} dua kali dengan
        \[
        \psi''(t) = \E_t(X^2) - \E_t(X)^2, \qquad
        \E_t(Y) := \frac{\E(Y\eu^{tX})}{\E(\eu^{tX})},
        \]
        yakni sebuah \emph{\omterm{def:b2:randomvar:variance}{ragam}} bagi peubah terbobot ulang yang masih
        bernilai di $\intcc01$; lalu batasilah ia dengan $\frac14$
        (lewat hujah keminimalan \cref{exo:b2:randomvar:6}) kemudian
        rampungkan lewat Taylor:
        \[
        \E\bigl(\eu^{t(X - p)}\bigr) \leq \eu^{t^2/8} .
        \]
  \item (Ketaksamaan Hoeffding) Untuk $X_i \in \intcc01$ yang \omterm{def:b2:proba:independence}{saling
        bebas} dengan rerataan bersama $p$, turunkanlah
        \[
        \P\bigl(\abs{\widehat p_n - p} \geq \delta\bigr)
        \leq 2\,\eu^{-2n\delta^2}
        \qquad (\delta > 0).
        \]
  \item Bandingkan laju Chebyshev
        $\frac{p(1-p)}{n\delta^2}$ dengan laju Hoeffding
        $2\eu^{-2n\delta^2}$: hipotesis mana yang dituntut masing-masingnya,
        dan mulai dari $n$ yang mana (secara kasar) batas
        eksponensialnya menang di $\delta = 0.03$ dan $p =
        \frac12$?
\end{enumerate}

\medskip
\noindent\textbf{Bagian III --- Berapa orang yang harus ditanya sebuah
jajak pendapat?} Sebuah jajak pendapat menanyai $n$ pemilih yang \omterm{def:b2:proba:independence}{saling bebas} dan dipilih seragam;
masing-masing menjawab jujur; dengan $p$ menyatakan skor benarnya dan $\widehat p_n$
angka jajak pendapatnya.
\begin{enumerate}[resume]
  \item Tunjukkan bahwa jajak pendapatnya akurat sampai $\pm\delta$ dengan
        kepercayaan $1 - \alpha$ (yakni
        $\P(\abs{\widehat p_n - p} \geq \delta) \leq \alpha$)
        segera setelah
        \[
        n \;\geq\; \frac{\ln(2/\alpha)}{2\,\delta^2} .
        \]
  \item Hitunglah $n$ yang dituntut bagi spesifikasi baku ``tiga
        poin, sembilan puluh lima persen'' (dengan $\delta
        = 0.03$ dan $\alpha = 0.05$): yakni $n \geq 2050$; dan bagi satu
        poin: $n \geq 18\,445$. Amatilah --- lalu jelaskanlah ---
        fakta yang mencolok bahwa jawabannya tak melibatkan
        besarnya populasi.
  \item Kerjakan ulang pertanyaan 13 lewat Chebyshev (dengan $V(X_1) = p(1-p)
        \leq \frac14$): yakni $n \geq \frac1{4\alpha\delta^2} =
        5556$ pada tiga poin. Perhatikan bahwa pencuplikan
        \emph{tanpa} pengembalian hanya membantu
        (\cref{exo:b2:randomvar:5}: sebab \omterm{def:b2:randomvar:variance}{ragamnya} menciut sebesar
        $\frac{N-n}{N-1}$).
  \item (Menetapkan pemenang pemilu) Skor benar seorang kandidat adalah $p
        = 0.52$. Berapa banyak pemilih yang harus dijajaki agar
        $\P(\widehat p_n \leq \tfrac12) \leq 0.01$? Tunjukkanlah $n
        \geq \frac{\ln 100}{2\cdot(0.02)^2} \approx 5757$ ---
        jadi menetapkan pemenang perlombaan yang ketat jauh lebih mahal daripada menaksir sebuah
        skor.
  \item Apa yang \emph{tak} diliput matematikanya: daftarkanlah
        asumsi pemodelan yang dipakai (yakni pencuplikan seragam yang \omterm{def:b2:proba:independence}{saling bebas},
        jawaban yang jujur, dan $p$ yang tetap), lalu jelaskanlah dalam
        satu paragraf pendek mengapa galat jajak pendapat yang nyata
        didominasi \emph{bias} (yakni pencuplikan yang tak seragam dan
        ketakresponsan), yang tak dikurangi oleh kenaikan $n$ mana pun.
\end{enumerate}

\medskip
\noindent\textbf{Bagian IV --- Lebih tajam dan lebih murah.}
\begin{enumerate}[resume]
  \item (Median rerata: peluruhan eksponensial dari dua momen)
        Pecahlah sebuah anggaran berisi $km$ cuplikan menjadi $k$ kelompok
        \omterm{def:b2:proba:independence}{saling bebas} berisi $m$; lalu misalkan $\widehat p^{(1)}, \dots, \widehat
        p^{(k)}$ rerata kelompoknya dan $M$ mediannya.
        Pilihlah $m$ sedemikian sehingga tiap kelompoknya memenuhi
        $\P(\abs{\widehat p^{(i)} - p} \geq \delta) \leq
        \frac18$ (lewat Chebyshev: sebab $m \geq \frac2{\delta^2}$
        sudah cukup). Tunjukkanlah bahwa bila $\abs{M - p} \geq \delta$
        maka sekurangnya $k/2$ kelompoknya keliru, lalu turunkan
        \[
        \P(\abs{M - p} \geq \delta)
        \leq \binom{k}{\lceil k/2\rceil}\Bigl(\frac18
        \Bigr)^{k/2}
        \leq 2^k\cdot 8^{-k/2} = 2^{-k/2} :
        \]
        jadi pemusatan eksponensial dengan memakai apa pun yang tak melampaui
        \omterm{def:b2:randomvar:variance}{ragamnya}.
  \item (Paley--Zygmund) Untuk $X \geq 0$ yang bermomen kedua,
        buktikanlah $\P(X > 0) \geq \dfrac{\E(X)^2}{\E(X^2)}$
        \emph{(lewat Cauchy--Schwarz pada $X\mathbf 1_{X>0}$)}: yakni
        perkakas arah baliknya --- sebab momen juga dapat memaksa
        \omterm{def:b2:proba:space}{kejadian} untuk terjadi.
  \item (Pinsker ringan) Tunjukkan $I(a) \geq 2\bigl(a -
        \tfrac12\bigr)^2$ pada $\intoo{\frac12}1$ \emph{(sebab
        selisihnya lenyap sampai orde dua di $\frac12$ dan
        turunan keduanya $\frac1{a(1-a)} - 4 \geq
        0$)}: jadi eksponen \omterm{def:b2:multint:exact}{eksak} Chernoff selalu mengalahkan
        eksponen kuadratik Hoeffding.
  \item Uraikan $I\bigl(\tfrac12 + \delta\bigr) = 2\delta^2 +
        O(\delta^4)$ lalu gabungkan dengan pertanyaan 4: jadi untuk
        simpangan yang kecil eksponen Hoeffding $2n\delta^2$ bersifat
        \emph{\omterm{def:b2:multint:exact}{eksak}} secara asimtotik --- sehingga tak ada metode yang dapat mengalahkannya
        lebih daripada faktor polinomial.
  \item Susunlah tabel kotak perkakasnya: yakni bagi Markov, Chebyshev, batas
        momen keempat pada \cref{exo:b2:randomvar:9},
        Hoeffding, dan Chernoff dengan eksponen $I$, nyatakanlah dalam
        satu baris masing-masing: hipotesis yang dituntut, peluruhan yang diperoleh,
        dan pertanyaan pada soal ini tempat ia paling tajam.
\end{enumerate}

\medskip
\noindent\textbf{Bagian V --- Panennya.}
\begin{enumerate}[resume]
  \item (Menguji sebuah koin) Sebuah koin entah setimbang entah berat sebelah dengan
        $p = 0.55$. Kamu melemparnya $n$ kali lalu mengumumkan
        ``berat sebelah'' ketika $\widehat p_n > 0.525$. Tunjukkanlah bahwa
        kedua peluang galatnya paling banyak
        $\eu^{-2n(0.025)^2}$, dan bahwa $n \geq 3685$ lemparan
        menjamin keduanya di bawah $1\%$.
  \item (\omterm{def:b2:proba:space}{Kejadian} langka menuntut batas yang sadar \omterm{def:b2:randomvar:variance}{ragam}) Misalkan $p =
        0.01$ lalu ambillah spesifikasi relatifnya $\delta =
        p/2 = 0.005$ dan $\alpha = 0.05$. Bandingkanlah ukuran
        cuplikan yang dituntut Hoeffding (yakni $n \approx 74\,000$) dan
        yang dituntut Chebyshev dengan \omterm{def:b2:randomvar:variance}{ragam} benarnya $p(1-p)$ (yakni $n
        \approx 7920$): jadi batas eksponensial yang buta \omterm{def:b2:randomvar:variance}{ragam}
        kalah oleh momen kedua yang sederhana. Nyatakanlah pelajarannya,
        dan dari mana perkakas yang hilangnya (yakni batas
        eksponensial yang sadar \omterm{def:b2:randomvar:variance}{ragam}; dan hampiran Poisson pada
        \cref{ch:b2:genfun}) kelak datang.
  \item (\omterm{def:b2:randomvar:law}{Hukum} kuat bagi koin) Dari $\sum_n
        2\eu^{-2n\delta^2} < \infty$ dan Borel--Cantelli
        (\cref{thm:b2:proba:borelcantelli}), buktikanlah bahwa
        $\widehat p_n \to p$ hampir pasti bagi lemparan koin yang \omterm{def:b2:proba:independence}{saling
        bebas}: yakni rumuskanlah \omterm{def:b2:proba:space}{kejadian} hampir pastinya sebagai
        $\bigcap_j\bigcup_N\bigcap_{n\geq N}
        \{\abs{\widehat p_n - p} < \tfrac1j\}$ seperti pada
        \cref{exo:b2:randomvar:9}, lalu simpulkan. (Adapun keterbatasannya
        menggantikan momen keempat yang dipakai di sana.)
  \item Rangkuman. Dalam lima kalimat: apa yang diharga dan dibeli tiap anak
        tangganya (yakni momen satu, dua, empat; eksponensial yang terbatas;
        dan eksponen yang \omterm{def:b2:multint:exact}{eksak}); mengapa menjajaki $2050$
        orang sudah cukup bagi sebuah negeri sebesar apa pun; dan
        batas mana di antaranya yang kelak dipertajam jilid Tahun ke-3 menjadi
        konstanta \omterm{def:b2:multint:exact}{eksak} teorema limit pusat.

\end{enumerate}
\end{problem}
